% Generated by report_gen.py (hwde v1) - do not hand-edit.
\documentclass[11pt]{article}
\newcommand\hspaper{a4paper}
\usepackage[nodiagramkit]{housestyle}
\usepackage{longtable}
\usepackage{pdfpages}
\makeatletter
\@ifundefined{KV@Gin@artifact}{\define@key{Gin}{artifact}[]{}}{}
\makeatother
\setcounter{tocdepth}{1}
\hsslug{PCB-0021-A \textperiodcentered\ lipo-boost}
\renewcommand{\sectionmark}[1]{\markboth{#1}{}}
\hssection{\leftmark}
\begin{document}
\sloppy
\hstitleblock{PCB-0021-A \textperiodcentered\ hwde v1 pipeline \textperiodcentered\ generated 2026-10-01 10:06}%
{lipo-boost \textemdash{} Highlights}%
{Phase P9; all 5 recorded gates pass.}


\section{At a Glance}
I want a little board that runs a 5 V USB gadget off a single-cell LiPo and charges the cell from USB.

\begin{itemize}
\item USB in (micro-B or C) charges a 1S LiPo on a JST-PH 2-pin connector, about 1 A max.
\item Load sharing: when USB is plugged in, the 5 V output keeps running while the cell charges.
\item 5 V / 1 A out on a USB-A socket, from a synchronous boost.
\item An enable pin to turn the output off, plus a low-battery LED and charge / done LEDs.
\item Small, cheap, JLCPCB-assembled, 2 layers is fine.
\end{itemize}
{\small
\begin{longtable}{lp{11cm}}
\hstoprule
\hshead{Field} & \hshead{Value} \\
\hline
\endhead
part number & PCB-0021-A \\
board & lipo-boost \\
phase & P9 (DFM) \\
size & 53.9 x 38.8 mm \\
parts & 33 placed parts, 24 BOM lines \\
gates & 5 of 5 recorded gates pass \\
\hline
\end{longtable}
}
\section{The Board}
\begin{center}
\includegraphics[width=0.48\textwidth]{reports/renders/PCB-0021-A_lipo-boost_top.png}\hfill
\includegraphics[width=0.48\textwidth]{reports/renders/PCB-0021-A_lipo-boost_bottom.png}\hfill
\\{\small top and bottom, as rendered from the board file}
\end{center}
\section{Copper Layers}
The outer copper layers as KiCad shows them, with every pad, track and zone labelled with its net (zoom in to read the small ones).
\clearpage
\subsection*{F.Cu}
\begin{center}
\includegraphics[width=\textwidth,height=0.82\textheight,keepaspectratio]{reports/layers/PCB-0021-A_lipo-boost_F_Cu.pdf}\\
{\small\texttt{reports/layers/PCB-0021-A\_lipo-boost\_F\_Cu.pdf}}
\end{center}
\clearpage
\subsection*{B.Cu}
\begin{center}
\includegraphics[width=\textwidth,height=0.82\textheight,keepaspectratio]{reports/layers/PCB-0021-A_lipo-boost_B_Cu.pdf}\\
{\small\texttt{reports/layers/PCB-0021-A\_lipo-boost\_B\_Cu.pdf}}
\end{center}
\subsection*{3D view}
\begin{center}
\includegraphics[width=0.8\textwidth]{reports/layers/PCB-0021-A_lipo-boost_iso.png}\\
{\small\texttt{reports/layers/PCB-0021-A\_lipo-boost\_iso.png}}
\end{center}
\section{What Is On It}
{\small
\begin{longtable}{p{3.6cm}p{5.4cm}p{3cm}l}
\hstoprule
\hshead{Refs} & \hshead{Part} & \hshead{Footprint} & \hshead{LCSC} \\
\hline
\endhead
C1 & 4.7uF 25V X5R & C0805 & C1779 \\
C2,C3,C10 & 10uF 25V X5R & C0805 & C15850 \\
C4,C13,C14 & 100nF 50V X7R & C0603 & C14663 \\
C11,C12 & 22uF 25V X5R & C0805 & C45783 \\
D1 & ORH-R36A orange & LED0603-RD & C205445 \\
D2 & KT-0603G green & LED0603-RD & C12624 \\
D3 & KT-0603R red & LED-SMD\_L1.6-W0.8-R-RD & C2286 \\
J1 & MicroXNJ & MICRO-USB-SMD\_MICROXNJ\_1 & C404969 \\
J2 & S2B-PH-SM4-TB(LF)(SN) & CONN-SMD\_P2.00\_S2B-PH-SM4-TB-LF-SN & C295747 \\
J3 & HX PZ2.54-1x7P & HDR-TH\_7P-P2.54-V-M & C32713273 \\
J4 & Type-A-AF90 & USB-A-TH\_U-A-24DD-Y-13 & C53207796 \\
L1 & 1uH & IND-SMD\_L4.0-W4.0 & C602028 \\
Q1 & 2N7002 & SOT-23-3\_L2.9-W1.3-P1.90-LS2.4-BR & C8545 \\
R1 & 10k & R0603 & C25804 \\
R2 & 1.2k & R0603 & C22765 \\
R3,R5,R7 & 1k & R0603 & C21190 \\
R4 & 51k & R0603 & C23196 \\
R6 & 470R & R0603 & C23179 \\
R8,R9,R11 & 100k & R0603 & C25803 \\
R10 & 732k & R0603 & C23239 \\
R12 & 0R & R0603 & C21189 \\
U1 & BQ24074RGTR & QFN-16\_L3.0-W3.0-P0.50-TL-EP1.7 & C54313 \\
U2 & TPS61023DRLR & SOT-563\_L1.6-W1.2-P0.50-LS1.6-BR & C919459 \\
U3 & MSXC61CN3302MR & SOT-23-3\_L2.9-W1.6-P1.90-LS2.8-BR & C45367187 \\
\hline
\end{longtable}
}
\includepdf[pages=1]{reports/schematic.pdf}
\section{Key Decisions}
{\small
\begin{longtable}{lp{1.6cm}p{6.2cm}p{6.2cm}}
\hstoprule
\hshead{Ph} & \hshead{When} & \hshead{Decision} & \hshead{Why} \\
\hline
\endhead
P2 & 10-01 05:51 & Closed both P2 research tasks --accept-drafts (6/12 verified, 6 refuted drafts stay draft and do not inject) & budget; refutations are numeric corrections carried as the decision below, not wrong rules \\
P2 & 10-01 05:51 & Design corrections from the second read: TPS61023 VREF 580/595/610 mV (PWM) for the FB divider tolerance; BQ24074 size ILIM/ISET with KILIM max 1720 (I = K/R); bq24074 VIN range 4.35-10.2 V (VOVP 10.2-10.8 V); TPS61023 RthJA 142.7 C/W -\textgreater{} \textasciitilde{}0.3-0.55 W loss is a 43-78 C rise on 2L, so give U2 a GND pour and vias; ISET pads need no dedicated vias & second reader refutations, research/records \\
P3 & 10-01 06:00 & PROVISIONAL: low-battery detector MSXC61CN3302MR C45367187 (3.3 V, not 3.2 V; releases \textasciitilde{}3.45 V) & no stocked 3.2 V open-drain detector; 3.3 V warns a little earlier, harmless \\
P3 & 10-01 06:00 & PROVISIONAL: RILIM 1.2k Basic (1.25-1.43 A) and RISET 1k Basic (\textasciitilde{}0.89 A) instead of 1.15k/910R Extended & two fewer extended-part fees; lower charge current is kinder to small cells; still about 1 A \\
P3 & 10-01 06:00 & bq24074 pins: no SYSOFF/TD; ITERM open (10 \%); EN2-\textgreater{}SYS, EN1 and CE-\textgreater{}GND (never float); TS 10k to GND; RTMR 51k (6.8 h) & part-sourcer read SLUS810N pin table; corrects architect's from-memory notes \\
P3 & 10-01 06:00 & U1 BQ24074RGTR stock thin (417), same-die alternates RGTT C2861252 / RGTRG4 & recheck before any order \\
P3 & 10-01 06:15 & J4 swapped C18077685 -\textgreater{} C53207796 Type-A-AF90 (JXTCONN USB-A THT right angle, stock 180) & C18077685 not in EasyEDA API: no pullable symbol/footprint \\
\hline
\end{longtable}
}
\section{How the Run Went}
22 recorded decisions from P2 to P7, 2026-10-01T05:43 to 2026-10-01T08:01. The run went back to an earlier phase 0 times.
\section{Versus Adafruit's PowerBoost 1000C}
This board is a from-scratch recreation of Adafruit's \textbf{PowerBoost 1000C} (designed by Limor Fried / Ladyada for Adafruit Industries; github.com/adafruit/Adafruit-PowerBoost-1000C), published under \textbf{CC BY-SA 3.0}. It was designed from the brief and datasheets, with no Adafruit file copied, so it is not a derivative; the original is credited as the reference.

\textbf{Same:} a load-sharing charger (about 1 A into a 1S LiPo, output stays up on USB), a 5 V / 1 A synchronous boost, an EN pin, low-battery, charge and done LEDs, a USB-A output and a JST-PH battery connector, on two copper layers.

\textbf{Different, and why:}
\begin{itemize}
\item Charger BQ24074 (QFN-16, about 1,066 in stock, \$1.23) rather than MCP73871 (QFN-20, about 112 in stock, \$2.79). Its fixed 4.4 V system rail keeps the boost boosting, never passing USB through to the output.
\item Boost TPS61023 (SOT-563, \$0.27) rather than TPS61090 (QFN-16, \$2.85). It has no low-battery comparator, so a separate 3.2 V detector (U3) drives the LBO output.
\item Output about 4.95 V, inside the USB-A limit; the original is about 5.2 V by design. Inductor 1 uH, not 6.8 uH.
\item Size 53.9 x 38.8 mm against about 36 x 23 mm for the original (its size is from the product page, not verified). Above the roughly 40 x 30 mm soft target in the requirements.
\item All 33 parts are JLCPCB-assembled; every IC is a JLC Extended part. A fixed 10 k resistor stands in for the cell thermistor.
\end{itemize}

\textbf{Open items for the owner} (provisional waivers, agent-proposed 2026-10-01; verify passes only with them):
\begin{itemize}
\item U1 thermal: 1.1 W worst case models about 81 C rise against 40 C declared. The BQ24074 folds charge current back at 125 C, so the cost is slower charging when hot.
\item U1 ground vias: 4 fit under the pad, the check wants 5.
\item +SYS pour neck and 2-via transition near R6: the DONE-LED branch, a few mA.
\item GND F.Cu neck beside USB-A J4: the return runs through its pads into the B.Cu plane; add stitching vias if F.Cu continuity is wanted.
\item Earlier provisional P2 choices: no input fuse or TVS, no reverse-battery FET, no 5 V limiter, 1 A default charge current.
\end{itemize}
Route style: 108 segments, 0 arcs, 0 off-angle, two small-jog warnings (a 0.6 mm sidestep on +5V and twelve sidesteps of 0.05 to 0.11 mm on LBO). The full design doc has the detail.

\section{Checks and Cost}
{\small
\begin{longtable}{llrr}
\hstoprule
\hshead{Gate} & \hshead{Status} & \hshead{Attempts} & \hshead{Failing} \\
\hline
\endhead
erc & pass & 3 & 0 \\
place & pass & 25 & 0 \\
drc\_routed & pass & 5 & 0 \\
verify & pass & 4 & 0 \\
dfm & pass & 3 & 0 \\
\hline
\end{longtable}
}
The full design doc (lipo-boost-design-full) has every decision, the renders and the run's history.

\end{document}
